%! Factoring Special Forms — Unit 1, Lesson 1.5 — Guided Notes (Answer Key)
\documentclass[10pt]{article}
\usepackage{atda-article}
\usepackage{atda-key}   % pulls in -boxes; do NOT also load -boxes
\newcommand{\TallMath}[1]{$\displaystyle #1\rule[-1.4em]{0pt}{3.2em}$}
\setlength{\workrowsep}{6pt}

\begin{document}

\pageheader{Unit 1, Lesson 1.5}{Guided Notes --- Answer Key}
% NAMESTRIP: no \namedateperiod here — mirrors the blank. Teacher-only prose goes
% in the lesson plan as \begin{teachernote}[Guided Notes], never in a key.

\vspace{0.15in}

\begin{objectivebox}
\small
By the end of this lesson, I will be able to\dots
\begin{enumerate}[leftmargin=1.5em, itemsep=2pt, topsep=3pt]
  \item recognize and factor a \textbf{difference of squares}, $a^2-b^2 = (a+b)(a-b)$, and explain
        why a \textbf{sum} of squares is prime over the integers;
  \item recognize and factor a \textbf{perfect square trinomial}, $a^2 \pm 2ab + b^2 =
        (a \pm b)^2$, by testing the middle term;
  \item factor a \textbf{sum or difference of cubes}, $a^3 \pm b^3 = (a \pm b)(a^2 \mp ab + b^2)$;
  \item factor \textbf{completely}, GCF first, and \textbf{verify a polynomial identity} by
        multiplying the factored form back out.
\end{enumerate}
\end{objectivebox}

\vspace{0.10in}

\boxguard
\begin{vocabbox}
\small
Fill in each term as we build it together.
\vocabans{Perfect square (of a term)}{A term that is some other term squared, such as $49$ ($=7^2$)
or $9x^2$ ($=(3x)^2$).}
\vocabans{Difference of squares}{A two-term polynomial $a^2-b^2$; it always factors as $(a+b)(a-b)$.
A \emph{sum} of squares, $a^2+b^2$, is prime over the integers.}
\vocabans{Perfect square trinomial}{A trinomial $a^2 \pm 2ab + b^2$ --- the square of a binomial;
it factors as $(a \pm b)^2$.}
\vocabans{Sum and difference of cubes}{$a^3+b^3 = (a+b)(a^2-ab+b^2)$ and $a^3-b^3 =
(a-b)(a^2+ab+b^2)$; the second factor is prime over the integers.}
\vocabans{Polynomial identity}{An equation between two polynomial expressions that is true for
every value of the variable; you verify one by multiplying out and comparing like terms.}
\end{vocabbox}

\vspace{0.10in}

\boxguard[12]
\begin{hookbox}
\small
\textbf{The lobby with a planter in the middle.} A school is retiling a \textbf{square} lobby that
measures $x$ feet on each side. A \textbf{square} planter, $4$ feet on each side, sits at the center
and is not tiled, so the tiled region is
\[ A(x) = x^2-16 \quad \text{square feet.} \]
The tile supplier will not quote an L-shaped region --- they quote \textbf{rectangles}, and they
want two dimensions.

Cut the tiled region into two pieces and slide one against the other: it fits a rectangle exactly.
What are that rectangle's two dimensions? (Check yourself at $x=10$, where the tiled area is $84$
square feet.)

\ansline{$x^2-16 = (x+4)(x-4)$, so the rectangle is $(x+4)$ ft by $(x-4)$ ft.}\\
\ansline{At $x=10$: $14$ ft by $6$ ft, and $14 \cdot 6 = 84 = 100-16$ square feet.}\\
\end{hookbox}

\vspace{0.10in}

\boxguard[24]
\begin{notesbox}{1. Difference of Squares --- $a^2-b^2 = (a+b)(a-b)$}
\small
Multiply $(a+b)(a-b)$ and watch the middle:
\[ (a+b)(a-b) = a^2 - ab + ab - b^2 = a^2-b^2. \]
The two middle terms are \emph{opposites}, so they cancel and a \textbf{two-term} polynomial is
left. Run that backwards and the factoring takes no search at all --- you read it off.

\textbf{The recognition test --- all three must hold:} exactly \textbf{two} terms; both terms are
\textbf{perfect squares}; the sign between them is a \textbf{minus}.

\renewcommand{\arraystretch}{1.9}
\begin{tabularx}{\linewidth}{@{}p{3.0cm} p{3.8cm} X@{}}
\rowcolor{mist}
\textbf{Binomial} & \textbf{Written as $a^2-b^2$} & \textbf{Factored form} \\
\hline
$x^2-36$    & \ans{$x^2-6^2$}      & \ans{$(x+6)(x-6)$} \\
$x^2-121$   & \ans{$x^2-11^2$}     & \ans{$(x+11)(x-11)$} \\
$9x^2-16$   & \ans{$(3x)^2-4^2$}   & \ans{$(3x+4)(3x-4)$} \\
$25x^2-81$  & \ans{$(5x)^2-9^2$}   & \ans{$(5x+9)(5x-9)$} \\
$x^2+49$    & \ans{not a difference} & \ans{prime} \\
\end{tabularx}

\vspace{6pt}
\textbf{The last row is the trap, and it is the same argument as Lesson 1.4.} A \emph{sum} of
squares gives the pattern no minus sign to use, and the product-and-sum search settles it: the only
integer pairs with product $49$ are
\begin{work}
  1 \cdot 49 &= 49, \quad \text{but } 1+49 = 50 \\
  7 \cdot 7 &= 49, \quad \text{but } 7+7 = 14 \\
  (-1)(-49) &= 49, \quad \text{but } -1+(-49) = -50 \\
  (-7)(-7) &= 49, \quad \text{but } -7+(-7) = -14
\end{work}
and none of them sums to $0$. So $x^2+49$ is \ans{prime} over the integers.
\end{notesbox}

\vspace{0.10in}

\boxguard[22]
\begin{notesbox}{2. Perfect Square Trinomials --- $a^2 \pm 2ab + b^2 = (a \pm b)^2$}
\small
This one also comes straight off a multiplication you already did --- warm-up item 1(b) squared a
binomial:
\[ (a+b)^2 = a^2+2ab+b^2 \qquad\text{and}\qquad (a-b)^2 = a^2-2ab+b^2. \]

\textbf{The recognition test --- all three must hold:} exactly \textbf{three} terms; the
\textbf{first and last} terms are perfect squares and both \textbf{positive}; the middle term is
exactly \ans{twice} the product of $a$ and $b$. The sign of the middle term is the sign that goes
inside the parentheses.

\renewcommand{\arraystretch}{1.9}
\begin{tabularx}{\linewidth}{@{}p{3.0cm} p{3.8cm} X@{}}
\rowcolor{mist}
\textbf{Trinomial} & \textbf{$a$, $b$, and $2ab$} & \textbf{Factored form} \\
\hline
$x^2+12x+36$    & \ans{$x$, $6$; \ $2ab=12x$}    & \ans{$(x+6)^2$} \\
$x^2-16x+64$    & \ans{$x$, $8$; \ $2ab=16x$}    & \ans{$(x-8)^2$} \\
$4x^2+20x+25$   & \ans{$2x$, $5$; \ $2ab=20x$}   & \ans{$(2x+5)^2$} \\
$x^2+8x+25$     & \ans{$x$, $5$; \ $2ab=10x \ne 8x$} & \ans{not a perfect square; prime} \\
\end{tabularx}

\vspace{6pt}
\textbf{Row 4 is why you test the middle term instead of trusting the ends.} Both $x^2$ and $25$ are
perfect squares, so the row \emph{looks} like a match --- but $2ab = 2(x)(5) = 10x$, not $8x$. Fall
back on Lesson 1.4: no integer pair has product $25$ and sum $8$, so this one is prime.
\end{notesbox}

\vspace{0.10in}

\boxguard[24]
\begin{notesbox}{3. Sum and Difference of Cubes --- the One Genuinely New Pattern}
\small
Two terms again, but now both are perfect \textbf{cubes}. Unlike a sum of squares, a \textbf{sum of
cubes does factor}:
\[ a^3+b^3 = (a+b)\left(a^2-ab+b^2\right) \qquad
   a^3-b^3 = (a-b)\left(a^2+ab+b^2\right) \]

\textbf{SOAP} keeps the signs straight: \textbf{S}ame sign as the original, \textbf{O}pposite sign
in the middle, \textbf{A}lways \textbf{P}ositive at the end.

\textbf{Why should you believe it?} Multiply the right-hand side out --- that is exactly what
``verify an identity'' means:
\begin{work}
  (a-b)\left(a^2+ab+b^2\right) &= a^3+a^2b+ab^2-a^2b-ab^2-b^3 \\
                               &= a^3-b^3
\end{work}
Everything except the two cubes cancels in pairs. A statement like this, true for
\ans{every} value of $a$ and $b$, is a \textbf{polynomial identity}, and multiplying out is how
you prove one.

\textbf{The packing block.} A solid foam cube $x$ inches on an edge has a $2$-inch cube cut out of
its center, so the foam left is $x^3-8$ cubic inches:
\begin{work}
  x^3-8 &= x^3-2^3 \\
        &= (x-2)\left(x^2+2x+4\right)
\end{work}
At $x=5$ that is $125-8 = 117$ cubic inches, and $(3)(25+10+4) = 3(39) = 117$ as well.

\textbf{Do not factor the second factor again.} It is \emph{not} a perfect square trinomial --- it
carries $ab$, not $2ab$ --- and $x^2+2x+4$ is prime over the integers.

\renewcommand{\arraystretch}{1.9}
\begin{tabularx}{\linewidth}{@{}p{3.0cm} p{3.8cm} X@{}}
\rowcolor{mist}
\textbf{Binomial} & \textbf{$a$ and $b$} & \textbf{Factored form} \\
\hline
$x^3+64$     & \ans{$a=x$, $b=4$}   & \ans{$(x+4)(x^2-4x+16)$} \\
$27x^3-1$    & \ans{$a=3x$, $b=1$}  & \ans{$(3x-1)(9x^2+3x+1)$} \\
$8x^3+125$   & \ans{$a=2x$, $b=5$}  & \ans{$(2x+5)(4x^2-10x+25)$} \\
\end{tabularx}
\end{notesbox}

\vspace{0.10in}

\boxguard[20]
\begin{notesbox}{4. Factoring Completely --- The Order, Finished}
\small
Lesson 1.3 started this list and Lesson 1.4 added a line. Today it is complete:
\begin{enumerate}[leftmargin=1.4em, itemsep=2pt, topsep=3pt]
  \item \textbf{GCF first, always.}
  \item \textbf{Count the terms.} \textbf{Two} $\Rightarrow$ difference of squares, or sum or
        difference of cubes (a sum of squares is prime). \textbf{Three} $\Rightarrow$ test for a
        perfect square trinomial first, then Lesson 1.4. \textbf{Four} $\Rightarrow$ group.
  \item \textbf{Check every factor again} --- a special form can hide inside another one --- then
        multiply back.
\end{enumerate}
\begin{work}
  2x^3-50x &= 2x\left(x^2-25\right) \\
           &= 2x(x+5)(x-5)
\end{work}
Step 3 finally has teeth. Here the first factor is finished, but the second is not:
\begin{work}
  x^4-16 &= \left(x^2+4\right)\left(x^2-4\right) \\
         &= \left(x^2+4\right)(x+2)(x-2)
\end{work}
$x^2-4$ was still a difference of squares; $x^2+4$ is a sum of squares, so it stops there.
\end{notesbox}

\vspace{0.10in}

\boxguard[20]
\begin{practicebox}
\small
\textbf{The second lobby.} The same crew takes a second job: a square lobby $2x$ feet on a side with
a square planter $7$ feet on a side at its center, so the tiled region is $A(x) = 4x^2-49$ square
feet.

\begin{enumerate}[leftmargin=1.6em, itemsep=6pt, topsep=4pt]
  \item Factor $A(x)$ to give the supplier the rectangle's two dimensions.
\begin{work}
  4x^2-49 &= (2x)^2-7^2 \\
          &= (2x+7)(2x-7)
\end{work}

  \item The lobby is built at $x=6$. Give both dimensions and the tiled area, in feet, with units
        --- then check the area against the original polynomial.

\ansline{$2x+7=19$ ft by $2x-7=5$ ft, so $19 \cdot 5 = 95$ sq ft, and $4(36)-49 = 95$ too.}

  \item A third job is quoted as $3x^3-27x$ square feet of tile. Factor it \textbf{completely}, then
        say what step 1 of the order saved you.
\begin{work}
  3x^3-27x &= 3x\left(x^2-9\right) \\
           &= 3x(x+3)(x-3)
\end{work}
\ansline{Pulling the $3x$ out first left $x^2-9$ --- a difference of squares you can read off.}\\
\ansline{Without it, $3x^3-27x$ has no special-form pattern at all.}\\
\end{enumerate}
\end{practicebox}

\end{document}
